\chapter*{Заключение} \label{ch-conclusion}
\addcontentsline{toc}{chapter}{Заключение}

В ходе выполнения практической работы был проведён статический и динамический
анализ приложения поиска палиндромных подстрок на Node.js. Для анализа были
выбраны два модуля с внесёнными ошибками: \texttt{datasetNameParser.js} и
\texttt{percentInputDecoder.js}.

Динамический анализ выполнен с помощью Jazzer.js. Фаззер обнаружил обе ошибки.
В модуле разбора имени датасета минимальным входом стала пустая строка, вызывающая
\texttt{TypeError} из-за обращения к \texttt{parts[2].toLowerCase()} без проверки
структуры имени. В модуле декодирования percent-encoded строки минимальным входом
стала строка \texttt{\%}, вызывающая \texttt{URIError} в \texttt{decodeURIComponent}.
Обе ошибки были найдены быстро: первая за 1 запуск, вторая за 134 запуска.

Статический анализ выполнен двумя наборами правил: ESLint Security и
SonarJS-совместимыми правилами ESLint. Анализаторы обнаружили предупреждения,
связанные с небезопасным регулярным выражением, динамическими путями файлов,
потенциальной object injection и неиспользуемой переменной. Для предупреждений
составлена таблица сопоставления с CWE и БДУ ФСТЭК России.

По результатам сравнения инструментов установлено, что фаззинг эффективен для
поиска конкретных аварийных входов и даёт точные crash-файлы. Статический анализ
эффективен для предварительного поиска классов дефектов без выполнения программы,
но требует ручной проверки ложноположительных предупреждений. Совместное
использование фаззинга и статического анализа даёт более полную картину качества
кода: фаззер подтверждает воспроизводимость ошибки, а анализатор показывает
потенциально опасные места до запуска приложения.
